\section{Independent Confirmation}
\label{sec:confirm}

\paragraph{Holdout.} The 34 telecom holdout tasks have not been run by any agent in this study. Telecom differs from retail and airline in its tools and requirements, and in dual control: the user acts on a simulated device, so both sides' tool calls change state. Before the holdout is opened, the mechanism categories found in development are mapped to telecom tools and requirements; the telecom development tasks may be used to debug the tool table and to check that the categories fit, with every rule change logged. The protocol, the tool-type table (by hash), the hypothesis, the number of blocks and the stopping rule are then frozen in a pre-registration file. The block count is derived from the development block variance, not fixed in advance.

\paragraph{Thresholds and their rationale.} The confirmation uses the thresholds of \S\ref{sec:inference}: $g_{\min}=3$\pp{} for $G$, and, for the background success comparison, a non-inferiority margin $\varepsilon=3$\pp{} (H$_0$: success$_{\mathrm{Q}}-$success$_{\mathrm{B}}\le-\varepsilon$; pass if the one-sided 95\% lower bound exceeds $-3$\pp{}). The margin was written down before any paired BF16 outcome existed. On the 85-task pool, 3\pp{} is a change of about 2.5 task outcomes per run. The run-level nonce SD of one checkpoint was 2.5\pp{} on q4 and 2.9\pp{} on q8, so a smaller margin could not be reached with a few paired blocks; a 5\pp{} margin would sit at the size of the largest W4 checkpoint effect we observed (q8 AWQ-W4 d3, $+5.3$\pp{} on average), so it could not flag such a checkpoint. Using one margin for $G$ and for success keeps ``negligible'' consistent. Non-inferiority is not equivalence: equivalence additionally needs the two one-sided tests, with the 90\% interval inside $(-3,+3)$\pp{}, and a margin the data cannot resolve is reported as inconclusive.

\paragraph{What counts as confirmation.} The frozen development hypothesis is tested on the holdout with the same estimands and annotation. With 34 tasks, few holdout events for a mechanism mean limited evidence for it, not its absence; and fixed-pool conclusions on the holdout do not by themselves generalise to new tasks (\S\ref{sec:inference}).
